% Generated draft: requires mathematical and source review.
\documentclass[11pt]{article}
\usepackage[margin=25mm]{geometry}
\usepackage{fontspec}
\setmainfont{Latin Modern Roman}
\usepackage{amsmath,amssymb,mathtools,unicode-math}
\setmathfont{Latin Modern Math}
\usepackage{enumitem,needspace,xurl,hyperref}
\setlength{\parindent}{0pt}
\setlength{\parskip}{5pt}
\setlength{\emergencystretch}{3em}
\newcommand{\N}{\mathbb{N}}
\newcommand{\Z}{\mathbb{Z}}
\newcommand{\Q}{\mathbb{Q}}
\newcommand{\R}{\mathbb{R}}
\newcommand{\C}{\mathbb{C}}
\newcommand{\F}{\mathbb{F}}
\newcommand{\A}{\mathbb{A}}
\newcommand{\PP}{\mathbb{P}}
\newcommand{\E}{\mathbb{E}}
\newcommand{\eps}{\varepsilon}
\newcommand{\dd}{\,\mathrm{d}}
\begin{document}
\section*{UnsolvedMath ID: 20000744}
\textbf{Source problem code:} catalogue entry
\subsection{Definitions and mathematical statement}
This record splices two independent workshop questions. The supplied verification identifies ``28.'' as page footer 2 followed by Problem 8, not one problem numbered 28. Both mathematical requests are retained under this permanent record.

Problem 7, attributed to A. Kantorovich, asks for small gaps among prime values \(a_n=\lfloor n\log n\rfloor\). The source reports \(\#\{p\le x:p=a_n\text{ for some }n\}\sim x/(\log x)^2\) and suggests analogous questions for Piatetski--Shapiro primes, without specifying their exponent range. It does not define ``small''. Assuming infinitely many such prime values, enumerate \(q_j=a_{n_j}\) increasingly and put \(h_j=n_{j+1}-n_j\). The report's exact scale dictionary is
\[
 \liminf_{j\to\infty}\frac{q_{j+1}-q_j}{\log q_j}
 =\liminf_{j\to\infty}h_j
\]
in the extended sense, and
\[
 \liminf_j\frac{q_{j+1}-q_j}{(\log q_j)^2}=0
 \quad\Longleftrightarrow\quad
 \liminf_j\frac{h_j}{\log n_j}=0.
\]
Since \(h_j\ge1\), the ordinary-prime GPY target of a zero liminf after division by \(\log q_j\) is obstructed by the mesh. A finite liminf on that scale is equivalent to bounded producing-index gaps; a zero liminf on the sequence's mean-gap \((\log q)^2\) scale is a distinct weaker target. Absolute bounded gaps are likewise incompatible with the increasing mesh. None of these targets is chosen silently as the unique meaning of the source question.

Problem 8, attributed to A. Granville, asks whether GPY, possibly assuming EH, can show infinitely many composite Mersenne numbers \(2^p-1\) with \(p\) prime. The exponent is reconstructed from the damaged superscript, not interpreted as the linear number \(2p-1\). C. Elsholtz mentions possible related work of M. R. Murty under Artin's primitive-root conjecture, without a citation or a stated theorem.

The report leaves this Mersenne request open. Its standard divisor restriction is that for odd prime \(p\), every prime divisor \(\ell\) of \(2^p-1\) has \(\operatorname{ord}_\ell(2)=p\), hence \(\ell\equiv1\pmod{2p}\); this is a restriction, not a construction of infinitely many composite cases. Its conditional Sophie-Germain route requires suitable primes \(p\equiv3\pmod4\) with \(2p+1\) prime and \(p>3\), so \(2p+1\) is a proper divisor of \(2^p-1\). Infinitely many such pairs would suffice, but that extra hypothesis is not proved or supplied by EH here. No GPY solution or verified Murty theorem is claimed.
\subsection{Short English statement}
The extracted record joins a thin-prime gap question and a separate prime-exponent Mersenne-compositeness question. The supplied report clarifies gap scales for \(\lfloor n\log n\rfloor\) and gives only restrictions/conditional routes for Mersenne composites. It does not solve either recovered request.
\subsection{Sources}
\begin{itemize}
\item Local UnsolvedMath registry, permanent ID \(20000744\), source code \texttt{catalogue entry}; AIM analytic number theory corpus, \texttt{aim-analytic-number-theory-notes.json}, notes[107], stored source item 7. Attribution: Ulam AI / UnsolvedMath Contributors, CC BY 4.0.
\item Angel Kumchev's notes for the December 2005 ARCC workshop \emph{Gaps Between Primes}, Problems 7 and 8: \url{https://aimath.org/WWN/primegaps/WorkshopProblems.pdf}. The supplied verification reports checking the five-page PDF and neighboring records and identifying the footer splice. This draft claims no new inspection.
\item The text names A. Kantorovich, A. Granville, C. Elsholtz and possible M. R. Murty work, with GPY, EH and Artin's conjecture; no full bibliographic entries or precise conditional theorem are supplied. Original remarks and literature fields are empty.
\item Registry dated August 14, 2026: statement \texttt{unrecoverable}, status open, classification \texttt{OPEN-TRIAGE}, suggested difficulty L2; attempt 1 labels the single spliced statement invalid and not a problem, rigor proved, novelty confidence low. The gap-scale reduction is a machine-generated research aid, not peer reviewed, requiring independent verification before citation. Both recovered questions are preserved despite that single-record diagnosis.
\end{itemize}
\end{document}
